\section{Related Work}\label{sec:related-work}

\subsection{Critic-Based and Critic-Free Language RL}

Modern language-model RL inherits its optimization machinery from policy-gradient reinforcement learning, from REINFORCE \citep{williams1992reinforce} to trust-region and clipped-surrogate methods such as TRPO, GAE, and PPO \citep{schulman2015trpo, schulman2016gae, schulman2017ppo}. In RLHF systems, these ideas are typically combined with learned value functions and reward models to stabilize policy optimization over text trajectories \citep{christiano2017preferences, ouyang2022instruct}. A recurring concern, however, is whether prefix-value estimation is the right abstraction for autoregressive language. Recent work revisits this question and shows that critic-free estimators can be surprisingly competitive. ReMax \citep{li2024remax}, Back to Basics \citep{ahmadian2024backtobasics}, and DeepSeekMath's GRPO formulation \citep{shao2024deepseekmath} all demonstrate that sequence-level baselines can recover much of the practical benefit of value-based methods without training a separate critic.

Our work is closest in spirit to this critic-free line: we keep the sequence-level objective and batch baseline, but change how the resulting scalar advantage is distributed over tokens. In other words, we study a credit-allocation layer that is orthogonal to whether the outer estimator is RLOO- or GRPO-style.

\subsection{Token-Level Credit Assignment and Reward Redistribution}

A large body of recent work argues that sequence-level rewards are too coarse and introduces denser token-level supervision. Some methods estimate token rewards with auxiliary models. PRIME \citep{cui2025prime} and related process-reward approaches learn implicit or explicit intermediate supervision for reasoning. Other methods redistribute a holistic reward to tokens algorithmically. RED \citep{huang2024red} derives token-level rewards from sequence scores, SCAR \citep{cao2025scar} uses Shapley-style attribution, VinePPO \citep{kazemnejad2024vineppo} estimates more accurate token credit via continuation rollouts, TEMPO \citep{tran2025tempo} uses tree structure, and GRPO-$\lambda$ \citep{parthasarathi2025grpolambda} propagates information backward with eligibility-trace-style weighting.

These papers make an important point: token-level credit assignment matters. They also clarify what our method is \emph{not}. We do not estimate token values, construct rollout trees, or train token reward models. Instead, we adopt a narrower objective: reweight the existing score-function terms using policy-internal salience signals already present in the forward pass.

\subsection{Intrinsic Token Salience and Selective Optimization}

The most relevant neighboring literature studies token importance using the policy itself. ConfPO \citep{yoon2025confpo} uses policy confidence to select critical tokens in preference optimization. Beyond the 80/20 Rule \citep{wang2025beyond8020} shows that a minority of high-entropy tokens drive much of RLVR improvement in reasoning models, suggesting that uniform token updates are unnecessarily diffuse. Token weighting has also proven useful outside RL: \citet{helm2025tokenweighting} show that non-uniform token emphasis can improve long-range language modeling in supervised settings.

These papers are the main novelty boundary for our work. We therefore make a deliberately narrow claim. We do not claim to be the first to identify non-uniform token salience, nor the first to use entropy- or confidence-like signals. Our specific contribution is to study \emph{continuous intrinsic token weighting for online verifiable-reward RL}: rather than hard-selecting a subset of tokens or estimating dense rewards, we continuously redistribute a sequence-level advantage with surprisal- and entropy-derived weights. This makes the method lightweight, differentiable through the outer objective, and easy to layer onto existing critic-free RLVR pipelines.
